\documentclass[../cours_1re.tex]{subfiles}
\begin{document}
\setcounter{chapter}{29}
\chapter{Relier les outils et construire le parcours vers le master}
\label{ch:30}
\begin{objectif}
Relier les outils construits et situer la progression vers un master.
Ce chapitre ne remplace pas les futurs cours de terminale, licence et master :
il montre des usages précis et propose des critères de maîtrise.
\end{objectif}
\section{Un même théorème dans trois langages}
En géométrie, si $p$ est le projeté de $x$ sur un sous-espace $F$,
$\|x-z\|^2=\|x-p\|^2+\|p-z\|^2$ pour $z\in F$.
En statistique, $\mathbb E((X-a)^2)=\operatorname{Var}(X)+(a-\mu)^2$.
Avec une information partielle, on a
$\mathbb E((X-Z)^2)=\mathbb E((X-Y)^2)+\mathbb E((Y-Z)^2)$.
Ces identités ont la même preuve : décomposer en deux termes orthogonaux,
développer un carré, annuler le terme croisé.
Savoir reconnaître cette structure est plus utile que mémoriser
trois recettes sans lien.

\section{Décomposer un système en directions simples}
Considérons la récurrence à deux composantes
\[
\begin{pmatrix}x_{n+1}\\y_{n+1}\end{pmatrix}
=
\begin{pmatrix}1/2&0\\0&1/4\end{pmatrix}
\begin{pmatrix}x_n\\y_n\end{pmatrix}.
\]
Elle donne $x_n=2^{-n}x_0$ et $y_n=4^{-n}y_0$ par deux suites géométriques.
Donc le vecteur tend vers $0$, et
$\sqrt{x_n^2+y_n^2}\le2^{-n}\sqrt{x_0^2+y_0^2}$.
Les axes sont des directions conservées par la transformation :
ce sont ici des directions propres.

Un exemple moins immédiat est la matrice
$A=\frac14\begin{pmatrix}3&1\\1&3\end{pmatrix}$.
On vérifie $A(1,1)=(1,1)$ et $A(1,-1)=\frac12(1,-1)$.
Décomposer $(x_0,y_0)=m(1,1)+d(1,-1)$, avec
$m=(x_0+y_0)/2$, $d=(x_0-y_0)/2$, donne
$A^n(x_0,y_0)=m(1,1)+2^{-n}d(1,-1)$.
Les deux composantes convergent vers leur moyenne commune.
Cette preuve annonce la diagonalisation sans supposer son théorème général.

\section{Du discret au continu}
La relation $u_{n+1}=qu_n$ mesure une multiplication par pas de temps.
L'équation $y'=ay$ mesure un taux instantané, avec $y(t)=y(0)e^{at}$.
Au pas $h$, son facteur exact est $e^{ah}$ ; Euler utilise $1+ah$.
L'approximation affine de l'exponentielle justifie leur proximité pour un petit pas.
Les matrices remplaceront ensuite les nombres pour des systèmes couplés.
Les séries du chapitre 26 suggèrent une exponentielle de matrice,
mais sa convergence et ses propriétés demanderont de définir une norme adaptée.

\section{Décomposer un signal discret}
Dans $\R^4$, les vecteurs
\[
v_0=(1,1,1,1),\quad v_1=(1,0,-1,0),\quad
v_2=(0,1,0,-1),\quad v_3=(1,-1,1,-1)
\]
sont deux à deux orthogonaux, comme on le vérifie en calculant leurs six
produits scalaires. Ils sont non nuls, donc libres ; tout vecteur de $\R^4$
se décompose sur eux. Plus explicitement, pour $x=(a,b,c,d)$,
\[
x=\frac{a+b+c+d}{4}v_0+\frac{a-c}{2}v_1
+\frac{b-d}{2}v_2+\frac{a-b+c-d}{4}v_3.
\]
Vérifier les quatre coordonnées démontre directement la formule.
Les vecteurs $v_1$ et $v_2$ sont des valeurs de cosinus et de sinus sur
quatre points régulièrement espacés du cercle. Cette petite décomposition
annonce la transformée de Fourier discrète.
Omettre une composante donne une projection sur les composantes conservées :
l'erreur au carré est la somme des carrés des composantes omises,
avec les facteurs de norme correspondants.

\section{La progression à poursuivre}
\begin{longtable}{>{\raggedright\arraybackslash}p{.18\linewidth}>{\raggedright\arraybackslash}p{.74\linewidth}}
\toprule Étape & Objets et capacités à construire\\\midrule
Terminale &
Limites et continuité systématiques, logarithme, intégration, probabilités
binomiales et concentration, géométrie dans l'espace. Avec mathématiques
expertes : nombres complexes, arithmétique, graphes et matrices.\\
Licence 1 &
Logique et rédaction autonome ; ensembles et applications ; analyse réelle,
suites et séries ; algèbre linéaire et systèmes ; calcul différentiel et intégral.\\
Licence 2 &
Espaces vectoriels abstraits, réduction des matrices, fonctions de plusieurs
variables, séries de fonctions, probabilités et méthodes numériques.\\
Licence 3 &
Topologie et espaces métriques, mesure et intégration, groupes et anneaux,
analyse complexe, probabilités, premiers choix de spécialisation.\\
Master &
Selon le parcours : analyse fonctionnelle, équations différentielles et aux
dérivées partielles, géométrie, algèbre, probabilités avancées, statistiques
et optimisation. Lire des preuves longues, résoudre des problèmes ouverts
à l'échelle du cours et produire un travail autonome.\\\bottomrule
\end{longtable}
Il n'existe pas un programme unique de master de mathématiques :
les spécialités diffèrent. Ce tableau est une progression pédagogique,
pas une équivalence de diplôme ni le programme d'une université donnée.

\begin{finchapitre}
\section{Comment savoir si une notion est maîtrisée}
Pour chaque outil, être capable de donner sa définition, de citer exactement
les hypothèses d'un résultat, de reconstruire sa preuve et de fournir un
contre-exemple si l'on supprime une hypothèse.
Savoir ensuite résoudre un problème où le nom de l'outil n'est pas annoncé,
expliquer les unités ou le modèle, puis vérifier une réponse par un second point de vue.

Les problèmes de synthèse demandent ce transfert : coût et optimisation,
croissances discrète et continue, projection et ajustement, conditionnement
et prédiction, approximation et erreur. Une rédaction doit faire apparaître
le choix de la méthode, pas seulement la suite des calculs.
Les compléments peuvent être relus après le programme de terminale :
ne pas tout comprendre au premier passage n'empêche pas de progresser.


\section{Utilité et prolongements}
Les liens entre domaines constituent un outil de travail : un carré
développé peut résoudre une question de géométrie, de statistique ou d'optimisation.

Le parcours vers le master se construira en ajoutant des objets et des preuves,
avec une autonomie croissante. Ce volume fournit des fondations de première
et des premiers ponts démontrés, pas la maîtrise prématurée de ces domaines.
\begin{retenir}
Définir, démontrer, chercher un contre-exemple, appliquer et relier.
Le niveau avancé se construit par la maîtrise progressive des hypothèses
et des raisonnements, autant que par l'accumulation de formules.
\end{retenir}
\end{finchapitre}
\end{document}
